\chapter{奇点解消}

本章节说明如何通过沿光滑中心连续爆破，把一个光滑复解析流形上的自反饱和抛物层的奇性解消掉。经过奇点解消，我们将得到一个局部阿贝尔抛物丛，其抛物度数在拉回下保持不变。

下述定理 \cite[Theorem 4.11]{grivaux2010chern} 起关键作用。

\begin{theorem}
设 \(\F\) 是复解析流形 \(X\) 上的一个解析凝聚层。则存在一个双亚纯态射
\[
\pi:\widetilde X\to X,
\]
它通过沿光滑中心连续爆破得到，使得 \(\F\) 的拉回模去挠层后是局部自由的，且挠子层支撑在例外除子上。
\end{theorem}

容易证明下述推论。

\begin{corollary}
设 \(\F\) 是局部自由层 \(\E\) 的一个子层，则存在一个双亚纯态射 \(\pi:\widetilde X\to X\)，使得 \(\pi^*(\F)\) 在 \(\pi^*(\E)\) 中的饱和化是 \(\pi^*(\E)\) 的一个子丛。
\end{corollary}

设 \(\F_*\) 是 \((X,D)\) 上的饱和自反抛物层，其中 \(X\) 是光滑复解析流形，\(D\) 是简单正规交叉除子。由定理 3.1，我们可把底层层 \(\F\) 解消成一个向量丛 \(\mathcal E\to \widetilde X\)，其中 \(\pi_1:\widetilde X\to X\) 是定理 4.1 所述的某个修改。为节省记号，即便之后继续爆破，我们仍记修改为
\[
\pi:\widetilde X\to X,
\]
而 \(\mathcal E\) 的拉回仍记作 \(\mathcal E\)。我们也约定，总是假定拉回层已经模掉挠子，并继续使用原记号表示拉回层。此外，我们也不改变在连续爆破过程中生成的例外除子不可约分支逆像的记号。\(D\)（及 \(D_i\)）的严格变换始终记作 \(D^*\)（及 \(D_i^*\)），例外除子记作 \(\E\)。不失一般性，总可假定 \(D^*\) 与 \(\E\) 横截相交。

接下来，为了获得局部阿贝尔性质，我们继续爆破，使得至多只有两个除子不可约分支的严格变换相交。限制到每个 \(D_i^*\) 上，我们有 \(\mathcal E|_{D_i^*}\) 上的一列子层滤链 \({}^i\F_a|_{D_i^*}\)。于是可以应用推论 4.2，使其成为 \(\mathcal E|_{D_i^*}\) 上的一列子丛滤链。类似地，通过进一步爆破，也可使
\[
{}^i\F_a|_{D_{ij}^*}\cap {}^j\F_b|_{D_{ij}^*}
\]
成为 \(\mathcal E|_{D_{ij}^*}\) 的一个子丛，其中 \(a,b\in[0,1)\)。于是，我们在 \(D^*\) 上得到了 \(\mathcal E\) 的一个相容抛物结构。记此局部阿贝尔抛物丛为 \(\mathcal E_*\)。

由引理 3.3、定义 3.11，可容易得到：对任意 Kähler 类 \(\omega\) 以及 \(k=1,2\)，有
\[
\ch_k(\mathcal E_*)\cdot \pi^*\omega=\ch_k(\F_*)\cdot \omega.
\]

因此有：

\begin{proposition}
设 \(\F_*\) 是一对 \((X,D)\) 上的饱和自反抛物层，其中 \(X\) 是光滑复解析流形，\(D\) 是简单正规交叉除子。则存在一个由沿光滑中心连续爆破得到的双亚纯态射
\[
\pi:\widetilde X\to X,
\]
使得拉回层 \(\pi^*(\F)\) 的无挠商 \(\mathcal E\) 在 \(D^*\) 上带有局部阿贝尔抛物结构。此外，对任意 Kähler 类 \(\omega\) 及 \(k=1,2\)，有
\[
\ch_k(\mathcal E_*)\cdot \pi^*\omega=\ch_k(\F_*)\cdot \omega.
\]
\end{proposition}